\documentclass[12pt, a4paper]{article}

%%------------------Preamble------------------%%

    %% 引入Package %%
\usepackage[margin=2cm]{geometry} % 上下左右距離邊緣3cm
\usepackage{fancyhdr}   % 頁首頁尾 fancy header
\usepackage{titling}    % 設定 title,author,date
\usepackage{mathtools,amsthm,amssymb} % 引入 AMS 數學環境
\usepackage[shortlabels,inline]{enumitem} %加強 enumerate, itemize 功能
\usepackage{setspace}
\usepackage{tcolorbox}
\usepackage{xcolor} % 用於表格上色
\usepackage{graphicx}      % 用於縮放表格寬度
\usepackage{array}
\usepackage{float} % for [H] float placement
\tcbuselibrary{breakable}
\usepackage{tabularx}
\usepackage{colortbl}

    %% 中文環境，僅限使用 XeLaTeX 編譯器 %%
\usepackage[CJKmath=true,AutoFakeBold=3,AutoFakeSlant=.2]{xeCJK}

\setCJKmainfont{Noto Serif CJK TC}
\setCJKsansfont{Noto Sans CJK TC}
\setCJKmonofont{Noto Sans Mono CJK TC}

\newCJKfontfamily\Kai{Noto Serif CJK TC}
\newCJKfontfamily\Hei{Noto Sans CJK TC}
\newCJKfontfamily\NewMing{Noto Serif CJK TC}

\XeTeXlinebreaklocale "zh"




\begin{document}

\title{T6 選擇題解答} % 設定標題
\author{李原廷} % 設定作者
\date{\today} % 設定完成日期，也可以使用 \date{\today}

\maketitle



\begin{enumerate}

    \item 甲去年年底購入一筆土地價值 50 萬元，估計土地出租每年可獲收益 1 萬元。若該筆土地每年應課徵 2‰ 地價稅，某甲於今年底將土地出售，租稅完全資本化，則該筆土地財產現值為多少元？
    \begin{enumerate}[label=\Alph*.]
        \item 50 萬元
        \item 45 萬元
        \item 40 萬元
        \item 30 萬元
    \end{enumerate}
    解答：
    假設市場利率為$r$，則：
    \[
    50\text{萬} = \frac{1\text{萬}}{r}
    \]
    \[
    r = \frac{1}{50} = 0.02 = 2\%
    \]
    地價稅：$10,000\times 2‰ = 1,000$，則可以得出稅額的折現值：
    \[
    T = \frac{\frac{1,000}{1.2}}{1-\frac{1}{1.2}} = \frac{1,000}{0.2} = 50,000
    \]
    售出時資產現值：$50\text{萬} -5\text{萬} = 45\text{萬}$

    另外一種算法：$V_t$為稅後土地現值
    \[
    V_t = \frac{Y-tV}{r}= \frac{10,000 - 0.002\times 500,000}{0.02} = 450,000
    \]
    其中：$t$為課徵稅率、$V$為原土地價值、$Y$為每年可獲得報酬

    \item 若需求曲線為負斜率，供給彈性等於零，對廠商課徵貨物稅，將如何影響產量與價格？
    \begin{enumerate}[label=\Alph*.]
        \item 產量減少，消費者支付的價格上升，廠商收到的價格下降
        \item 產量減少，消費者支付的價格上升，廠商收到的價格不變
        \item 產量不變，消費者支付的價格不變，廠商收到的價格下降
        \item 產量不變，消費者支付的價格上升，廠商收到的價格不變
    \end{enumerate}
    解答：
    \begin{enumerate}[1.]
        \item 供給彈性等於零的意義：當供給彈性等於零（\(E_s = 0\)），代表供給曲線是一條垂直於數量軸的直線（完全無彈性）。這意味著無論市場價格如何變動，廠商所提供的產量都是固定不變的。
        \item 對產量的影響：因為供給量被完全鎖定在垂直線的位置上，所以即使政府對廠商課徵貨物稅，新的均衡數量依然會停留在該垂直線與原需求曲線的交點。因此，「產量不變」。
        \item 對消費者支付價格的影響：市場上的實際交易產量既然沒有改變，且需求曲線（消費者願意支付的價格軌跡）也沒有移動，所以消費者為了這個固定產量所願意支付的價格，會和課稅前完全一樣。因此，「消費者支付的價格不變」。
        \item 對廠商收到價格的影響：在市場中，消費者支付的價格不變，但政府從中抽走了貨物稅（\(t\)）。廠商實際收到的價格公式為：\(\text{廠商收到的價格} = \text{消費者支付的價格} - \text{稅額}\)
    \end{enumerate}
    既然消費者支付的錢沒變多，政府又要拿走一部分，廠商能收到的錢自然就變少了。因此，「廠商收到的價格下降」（以價格下降的形式呈現）。

    \item 對廠商所製造之產品課稅後，其轉嫁程度與供給彈性大小之關係，在下列之敘述中何者正確？
    \begin{enumerate}[label=\Alph*.]
        \item 前轉與產品供給彈性成反比，而後轉與要素供給彈性亦成反比
        \item 前轉與產品供給彈性成正比，而後轉與要素供給彈性亦成正比
        \item 前轉與產品供給彈性成正比，但後轉與要素供給彈性卻成反比
        \item 前轉與產品供給彈性成反比，但後轉與要素供給彈性卻成正比
    \end{enumerate}
    解答：
    租稅轉嫁：旨法定租稅與經濟歸宿之間差額。
    \begin{itemize}
        \item 法定歸宿：誰是法律上，被要求繳納該睡的人。
        \item 經濟歸宿：課稅造成私人實質負擔所得分配之改變。表示真正的最後實際負擔租稅者。
        \item 前轉：納稅人將租稅負擔轉嫁給買方，納稅人透過提高財貨價格的方法，將租稅轉嫁給消費者。
        \item 後轉：納稅人將租稅負擔轉嫁給賣方，納稅人透過降低財貨價格的方法，將租稅轉嫁給供給者。
    \end{itemize}
    廠商的財貨若供給彈性越大，越容易將稅負轉嫁給消費者，所以前轉與供給彈性成正比；若要素供給彈性越大（勞動供給曲線），越容易將稅負轉嫁給生產者，要素的供給者（勞動）承擔的負擔越少，所以後轉與要素供給彈性成反比。

    \begin{itemize}
        \item 前轉（Forward Shifting）與「產品」供給彈性的關係：
        \begin{itemize}
            \item 根據租稅負擔公式，消費者承擔的比例為：\(\frac{E_s}{E_d + E_s}\) （其中 \(E_s\) 為產品供給彈性，\(E_d\) 為產品需求彈性）。
            \item 當產品的供給彈性 (\(E_s\)) 越大時，表示廠商調整產量或退出市場的能力越強。如果消費者不願意承擔較高的含稅價格，廠商大可減少生產。因此，廠商越有能力將租稅往前轉嫁給消費者。
        \end{itemize}
        \item 後轉（Backward Shifting）與「要素」供給彈性的關係：
        \begin{itemize}
            \item 後轉是指廠商將租稅負擔藉由「壓低生產要素價格（如扣減工資、壓低原料進價）」的方式，轉嫁給要素供給者（勞工或原料商）。
            \item 此時廠商是買方，要素提供者是賣方。要素提供者承擔負擔的比例取決於他們自己的「要素供給彈性」。
            \item 當要素的供給彈性越大時，表示要素提供者（如勞工）越容易找到其他工作或將資源轉移到其他產業。如果廠商試圖壓低工資來轉嫁租稅，這些勞工就會直接離開。因此，廠商就越不可能將租稅往後轉嫁給要素提供者。
        \end{itemize}
    \end{itemize}

    \item 假設一經濟體系只有兩財貨：商品 \(X\)、商品 \(Y\) 與兩要素：勞動 \(L\)、資本 \(K\)，而 \(X\) 為資本密集性商品，\(Y\) 為勞動密集性商品，倘若政府只對 \(X\) 商品課稅，則下列有關其經濟效果之敘述，何者正確？
    \begin{enumerate}[label=\Alph*.]
        \item 當 \(X\) 商品之需求彈性愈大時，將導致 \(X\) 商品之價格上漲的幅度愈大
        \item 當 \(X\) 商品之需求彈性愈大時，將導致資本報酬率下降的幅度愈大
        \item 當 \(X\) 商品與 \(Y\) 商品之要素比例差異程度愈大時，將導致資本報酬率下降的幅度愈小
        \item 當 \(Y\) 商品的生產過程中，其資本愈難取代勞動時，將導致資本報酬率下降的幅度愈小
    \end{enumerate}
    解答：
    \begin{enumerate}[label=\Alph*.]
        \item 錯誤：前轉（價格上漲轉嫁給消費者）與需求彈性成反比。若 X 的需求彈性愈大，廠商越不敢漲價（否則消費者會大量流失），因此 X 價格上漲的幅度會愈小，租稅負擔會更多地往後轉嫁給生產要素。
        \item 正確：當 X 商品的需求彈性愈大時，代表消費者對含稅價格的反應越劇烈。課稅會導致 X 的產量大幅萎縮。產量萎縮得越多，釋放出的閒置資本就越多，市場上資本過剩的情況就越嚴重。因此，為了讓這些龐大的多餘資本被吸收，資本報酬率必須下降更大的幅度。
        \item 錯誤：當兩部門的要素密集度（資本/勞動比例）差異愈大時，代表 Y 產業極度偏好勞動，非常不喜歡用資本。當 X 釋出大量資本時，Y 產業非常難以消化這些資本，這會導致資本報酬率必須下降愈大的幅度，才能勉強誘使企業多使用資本。
        \item 錯誤：當 Y 產業的要素替代彈性越小（資本愈難取代勞動），代表就算資本變得再便宜，Y 產業的生產技術也很難改成多用資本。既然 Y 產業無法輕易透過改變生產技術來吸收市場上過剩的資本，那麼資本報酬率就必須被迫下跌愈大的幅度，才能達到市場均衡。
    \end{enumerate}

    \item 假設消費者的效用函數與預算限制式分別為 \(U(X, Y) = \sqrt{XY}\) 與 \(P_X X + P_Y Y = I\)，若 \(P_X = P_Y = 1\)，\(I = 100\)，下列三種稅制何種對消費者產生無謂損失最大？稅制甲：對消費者課徵 20\% 的所得稅；稅制乙：對消費者課徵 10\% 的一般消費稅；稅制丙：對消費者消費的 \(X\) 財貨課徵 5\% 的特種消費稅。
    \begin{enumerate}[label=\Alph*.]
        \item 稅制甲
        \item 稅制乙
        \item 稅制丙
        \item 三者的無謂損失相同
    \end{enumerate}
    解答：
    \begin{enumerate}[label=\Alph*.]
        \item 稅制甲：原預算式改寫為$\frac{P_X}{1-0.2} X + \frac{P_Y}{1-0.2} Y = I$沒有改變相對價格，不會有超額負擔。
        \item 稅制乙：原預算式改寫為$(1+0.1)P_X X + (1+0.1)P_Y Y = I$沒有改變相對價格，不會有超額負擔。
        \item 稅制丙：原預算式改寫為$(1+0.05)P_X X + P_Y Y = I$有改變相對價格，會有超額負擔。
    \end{enumerate}

    \newpage

    \item 在兩部門兩生產要素的 Harberger 一般均衡租稅歸宿模型中，若政府只對使用在其中一部門的勞動課稅，且勞動與資本可以替代，則下列敘述何者正確？
    \begin{enumerate}[label=\Alph*.]
        \item 資本供給者將負擔全部的租稅
        \item 若被課稅部門相對另一部門而言較勞動密集，則稅後資本報酬率對工資率的比必定下降
        \item 若被課稅部門相對另一部門而言較勞動密集，則稅後資本報酬率對工資率的比必定上升
        \item 若被課稅部門相對另一部門而言較資本密集，則稅後資本報酬率對工資率的比必定上升
    \end{enumerate}
    解答：略。

    \item 若課徵土地稅產生完全的資本化（capitalization）現象，下列關於課徵土地稅之敘述，何者正確？
    \begin{enumerate}[label=\Alph*.]
        \item 不會產生超額負擔
        \item 買方與賣方各負擔一部分租稅
        \item 土地成交量將減少
        \item 政府收到的稅收為零
    \end{enumerate}
    解答：
    \begin{enumerate}[1.]
        \item 土地供給完全無彈性：在經濟學中，土地的總數量是大自然固定的，因此一般假設土地的供給曲線為垂直線（供給彈性 \(E_s = 0\)）。
        \item 不會產生超額負擔（選項 A 正確，選項 C 錯誤）：超額負擔（無謂損失）的產生是由於租稅扭曲了市場價格，導致交易「數量」減少而產生的福利損失（公式為 \(DWL = \frac{1}{2} \cdot t \cdot \Delta Q\)）。因為土地供給彈性為零，課稅後土地的供給量與成交量並不會減少（\(\Delta Q = 0\)），因此不會產生超額負擔，這符合租稅中立性。
        \item 租稅完全資本化與負擔歸宿（選項 B 錯誤）：「租稅完全資本化」意指未來所有應繳的地價稅，都會折現並反映在當期土地價格的下跌中。因為土地供給完全無彈性，所以未來的買方會要求將未來的稅負從買價中完全扣除。這導致現在的土地所有權人（賣方）承受了 100\% 的租稅負擔（以財產價值減損的形式出現），買方實質上並未負擔任何租稅。因此，買賣雙方並非各負擔一部分。
        \item 政府稅收（選項 D 錯誤）：政府依然會每年向土地所有權人收取土地稅，因此稅收絕對不為零。租稅資本化只是解釋了這筆稅收「實質上是由誰買單」（即課稅當下的地主），並不影響政府徵收稅款的事實。
    \end{enumerate}

    \item 在封閉經濟體的兩部門模型分析架構下，兩部門以 \(L\) 及 \(K\) 兩要素，分別生產財貨 \(F\) 及財貨 \(M\)。在沒有儲蓄、稅率相同情況下，以下租稅歸宿等價關係（tax equivalence relations）之討論，何者正確？
    \begin{enumerate}[label=\Alph*.]
        \item 對 \(F\) 部門的兩要素課稅，等同對 \(M\) 部門之產出課稅
        \item 對 \(F\) 及 \(M\) 兩部門所使用的要素 \(L\) 課稅，等同對兩部門產出課稅
        \item 對兩產出 \(F\) 及 \(M\) 課稅，等同對兩要素 \(L\) 及 \(K\) 課稅
        \item 對於 \(F\) 部門之要素 \(L\) 及 \(M\) 部門之要素 \(K\) 課稅，等同課徵一般所得稅
    \end{enumerate}
    解答：略，表格。

    \item 假設某一土地每期期末之報酬為 \$100。每期期末課以地價稅稅額 \$2，在租稅完全資本化的假設下，若折現率為 10\%，下列敘述何者正確？1.該土地稅前收益價格為 \$1,100 2.該土地稅後收益價格為 \$1,080 3.租稅資本化的數額為 \$20 4.該土地稅後收益價格為 \$980
    \begin{enumerate}[label=\Alph*.]
        \item 1.,2.,3.
        \item 1.,3.,4.
        \item 僅2.,3.
        \item 僅3.,4.
    \end{enumerate}
    解答：
    \begin{enumerate}[1.]
        \item 稅前收益：$\frac{100}{0.1} = 1,000$
        \item 稅後收益：$\frac{100-2}{0.1} = 980$
        \item 租稅折現值：$\frac{2}{0.1} = 20$
    \end{enumerate}

    \item 在哈柏格（A. C. Harberger）之兩部門-兩要素一般均衡模型架構下，對於其中一部門之產出課徵貨物稅，部分租稅負擔可後轉，由該課稅貨物生產過程使用相對密集的要素負擔。下列有關後轉幅度的討論，何者錯誤？
    \begin{enumerate}[label=\Alph*.]
        \item 與課稅貨物之需求彈性大小成正比
        \item 與未課稅部門要素替代容易度成反比
        \item 與兩部門要素密集度的差異成反比
        \item 與稅率成正比
    \end{enumerate}
    解答：假設對 X 部門課徵貨物稅：
    \begin{enumerate}[label=\Alph*.]
        \item 正確：如果 X 的需求彈性很大，課稅會導致 X 的產量大幅度萎縮。產量萎縮得越嚴重，釋放出的要素就越多，市場上的超額供給就越大，導致該要素報酬率必須下降越大的幅度才能結清市場。
        \item 正確：未課稅部門（Y 產業）必須負責吸收這些多出來的要素。如果 Y 產業的要素替代彈性很大（很容易用資本取代勞動），它就能輕易吸收這些釋出的資本，資本報酬率只需微幅下降即可。反之，若替代困難（彈性小），資本報酬率就必須大跌才能強迫 Y 產業使用。因此，兩者呈反比。
        \item 錯誤：當兩部門的要素比例（密集度）差異越大時，代表釋出要素的比例與吸收要素的比例不平均。這種的錯位會導致市場難消化多餘的要素，因此該要素的報酬率必須出現更大的跌幅。所以後轉幅度與要素密集度的差異應該是成正比，而非反比。
        \item 正確：課徵的稅率越高，對價格與產量的扭曲就越嚴重，釋出的要素數量就越多，自然會導致要素報酬率下降的幅度越大。兩者呈正比關係。
    \end{enumerate}

    \item 下列何者是租稅資本化（tax capitalization）的現象？
    \begin{enumerate}[label=\Alph*.]
        \item 賣方將需繳交的稅負，藉由提升生產效率以降低生產成本的方式減少租稅的實際負擔
        \item 買方考慮購買的土地在未來需負擔地價稅負，因而降低土地的買價
        \item 賣方將需繳交的稅負，藉由提高價格的方式移轉給買方
        \item 賣方將需繳交的稅負，藉由提高賣給同業的價格的方式移轉給同業
    \end{enumerate}
    解答：租稅資本化是指租稅流量被包含在資產價格內的過程，課稅標的需為耐久性與收益性財產。由於租稅資本化，所有現在和未來的土地稅負將由賣方所承擔，屬於後轉形式。

    tax capitalization: Special kind of backward shifting a stream of taxes becomes incorporated into the price of an asset.

    \item 依據哈柏格對部分要素稅一般均衡分析之說明，若僅對公司部門之資本要素報酬課稅，引起資本要素流向非公司部門之經濟行為改變，下列敘述何者錯誤？
    \begin{enumerate}[label=\Alph*.]
        \item 係由公司部門與非公司部門之資本報酬差異引發經濟行為改變
        \item 要素替代效果使非公司部門資本使用量增加
        \item 產量效果使公司部門資本使用量減少
        \item 最終效果將使非公司部門之資本稅後邊際報酬高於公司部門
    \end{enumerate}
    解答：
    \begin{enumerate}[label=\Alph*.]
        \item 正確：課稅之初，公司部門的資本稅後報酬率會瞬間下降，低於非公司部門。資本擁有者為了追求更高的利潤，會將資本從稅後報酬較低的公司部門，轉移到非公司部門。因此，資本報酬的差異是引發資本流動的原因。
        \item 正確：由於對公司部門的資本課稅，使得公司部門使用資本的相對成本變貴。因此，公司部門會改變生產技術，少用資本、多用勞動。釋放出來的資本流向非公司部門，導致非公司部門的資本變得相對便宜，非公司部門就會多用資本、少用勞動。因此，要素替代效果會使非公司部門的資本使用量增加。
        \item 正確：公司部門因為被課稅，生產成本增加，導致其產品價格上升。消費者因此減少購買公司部門的產品，轉而購買非公司部門的產品。公司部門的產量減少，連帶使得對資本（及勞動）的需求減少，因此產量效果會使公司部門的資本使用量減少。
        \item 錯誤：在一般均衡分析中，資本會持續在兩個部門之間流動，直到無利可圖為止。也就是說，資本會不斷從公司部門流向非公司部門，這會使得非公司部門的資本供給增加（報酬率逐漸被稀釋下降），而公司部門的資本供給減少（報酬率逐漸回升）。最終市場達到均衡的條件是：兩部門的「稅後」資本報酬率必定相等（若不相等，資本就會繼續流動）。因此，最終結果並不會使非公司部門的稅後報酬高於公司部門，而是兩者相等。
    \end{enumerate}

    \item 若某獨占市場的財貨之需求曲線為負斜率直線，邊際成本線為正斜率直線，政府對其課徵每單位 \(t\) 元的從量稅，下列敘述何者錯誤？
    \begin{enumerate}[label=\Alph*.]
        \item 獨占廠商面對之稅前與稅後平均收入曲線垂直差距為政府每單位從量稅金額
        \item 獨占廠商面對之稅前與稅後邊際收入曲線為平行左移
        \item 稅後均衡交易數量較課稅前減少
        \item 由於市場獨占，廠商可將從量稅完全轉嫁給消費者負擔
    \end{enumerate}
    解答：
    假設需求曲線為線性負斜率：\(P = a - bQ\)

    邊際收益曲線則為：\(MR = a - 2bQ\)

    假設邊際成本線為線性正斜率：\(MC = c + dQ\)
    \begin{itemize}
        \item 課稅前均衡： \(MR = MC \Rightarrow a - 2bQ = c + dQ\)，解得 \(Q_0 = \frac{a - c}{2b + d}\)
        \item 對廠商課徵單位稅 \(t\)： 新的邊際成本變為 \(MC_t = c + dQ + t\)
        \item 課稅後均衡： \(MR = MC_t \Rightarrow a - 2bQ = c + dQ + t\)，解得新產量 \(Q_1 = \frac{a - c - t}{2b + d}\)
        \item 價格變化（轉嫁程度）：
    將產量代回需求函數，可算出價格上漲的幅度：
    \[
    \Delta P = P_1 - P_0 = b(Q_0 - Q_1) = b \cdot \left(\frac{t}{2b + d}\right) = t \cdot \left(\frac{b}{2b + d}\right)
    \]
    \end{itemize}
	結論：因為需求斜率絕對值 $b > 0$ 且 MC 斜率 $d > 0$，所以$ \frac{b}{2b + d}$必然小於 1。這表示價格上漲的幅度小於稅額 $t$ ，獨占廠商必須自行吸收一部分的租稅，無法完全轉嫁給消費者。

    \item 兩期消費模型中，在相同稅收下，分別考慮利息收入課稅、利息支出可當費用扣除的比例所得稅，及對兩期消費課徵相同稅率的比例消費稅，對家計單位消費決策的影響，在不考慮休閒的假設下，下列關於兩種課稅方式之經濟效果的敘述，何者正確？
    \begin{enumerate}[label=\Alph*.]
        \item 課徵所得稅會使第一期消費相對價格上升
        \item 課徵消費稅會使第一期消費相對價格下降
        \item 課徵所得稅會產生超額負擔
        \item 課徵消費稅會產生超額負擔
    \end{enumerate}
    解答：
    在不考慮休閒（即勞動供給固定）的兩期消費模型中，我們探討的是家計單位如何在「第一期消費 (\(C_1\))」與「第二期消費 (\(C_2\))」之間做跨期選擇。兩者的相對價格（即放棄一單位 \(C_1\) 可以換取多少 \(C_2\)）是由市場利率 (\(r\)) 決定的，相對價格為 \(1 + r\)。
    \begin{itemize}
        \item 課徵「比例消費稅」（ B 錯誤、 D 錯誤）：如果政府對兩期的消費課徵相同稅率 (\(t_c\)) 的消費稅。
        \begin{itemize}
            \item 第一期消費的代價變成 \((1+t_c)C_1\)
            \item 第二期消費的代價變成 \((1+t_c)C_2\)
            \item 跨期相對價格為：\(\frac{(1+t_c)(1+r)}{(1+t_c)} = 1+r\)
            \item 結論： 因為兩期稅率相同，相對價格 \(1+r\) 並未改變。這代表消費稅沒有扭曲消費者跨期的選擇，它發揮了類似「定額稅」的功能，因此不會產生超額負擔（無謂損失）。所以 B 和 D 是錯的。
        \end{itemize}
        \item 課徵「比例所得稅，且利息課稅/扣除」（ A 錯誤、 C 正確）：如果政府課徵所得稅 (\(t_y\))，因為所得稅會對「利息收入課稅」，也會讓「利息支出抵稅」，這會導致家計單位面對的稅後淨利率下降。
        \begin{itemize}
            \item 稅後淨利率變為：\(r' = r(1-t_y)\)
            \item 第一期消費的相對價格（機會成本）由 \(1+r\) 變成了 \(1 + r(1-t_y)\)。
            \item 因為 \(1 + r(1-t_y) < 1+r\)，這代表第一期消費的相對價格下降了（變得相對便宜），而不是上升，所以 A 錯誤。
            \item 結論： 因為所得稅改變了兩期消費的相對價格（扭曲了價格機能），促使人們減少儲蓄、增加當期消費，這種人為的扭曲就會產生超額負擔。所以 C 是正確的。
        \end{itemize}
    \end{itemize}

    \newpage

    \item 下列關於個人所得稅課徵對勞動供給影響之敘述，何者正確？
    \begin{enumerate}[label=\Alph*.]
        \item 若休閒為正常財，則替代效果增加勞動供給，而所得效果降低勞動供給
        \item 若休閒為正常財，則替代效果與所得效果皆降低勞動供給
        \item 若休閒為劣等財，則替代效果降低勞動供給，而所得效果增加勞動供給
        \item 若休閒為劣等財，則替代效果與所得效果皆降低勞動供給
    \end{enumerate}
    解答：政府課徵個人所得稅，會導致勞工的稅後實質工資率下降。工資率可以視為「休閒的機會成本（相對價格）」。透過「替代效果」與「所得效果」來分析工資下降對勞動供給（或休閒時間）的影響：
	\begin{enumerate}[1.]
	    \item 替代效果（Substitution Effect）：
        \begin{itemize}
            \item 定義：假設維持原有效用（實質購買力不變），單純因為相對價格改變所引起的消費變化。
            \item 分析：課稅使得稅後工資下降，代表「休閒」變得相對便宜了。因此，人們會選擇「多休閒、少工作」。
            \item 結論：無論休閒是正常財還是劣等財，替代效果必然會降低勞動供給。
        \end{itemize}
        \item 所得效果（Income Effect）：
        \begin{itemize}
            \item 定義：相對價格不變下，單純因為實質所得（購買力）變動所引起的消費變化。
            \item 分析：課稅使得勞工的實質所得（購買力）變少，變得更窮了。
            \item 情境 I（若休閒為正常財）：當人變窮時，對正常財的消費會減少。因此勞工會減少休閒，為了彌補被扣掉的稅而多去加班賺錢。所以，若休閒為正常財，所得效果會增加勞動供給。（選項 A、B 敘述皆錯誤）
            \item 情境 II（若休閒為劣等財）： 當人變窮時，反而會增加對劣等財的消費。因此勞工會增加休閒時間。所以，若休閒為劣等財，所得效果會降低勞動供給。
        \end{itemize}
	\end{enumerate}

    \item 假定利息所得要課稅但利息支出不可以扣抵，下列敘述何者正確？
    \begin{enumerate}[label=\Alph*.]
        \item 借款者將增加借款：儲蓄者將增加儲蓄
        \item 借款者的情況不變：儲蓄者儲蓄的增減，視所得效果和替代效果的相對大小而定
        \item 借款者將減少借款：儲蓄者將減少儲蓄
        \item 儲蓄者的情況不變：借款者借款的增減，視所得效果和替代效果的相對大小而定
    \end{enumerate}
    解答：
    將市場上的消費者分為兩類：「儲蓄者」（第一期消費小於所得）與「借款者」（第一期消費大於所得），並分析該租稅政策對兩者「跨期預算線」的影響。
    \begin{enumerate}[1.]
        \item 對「借款者」的影響：
        \begin{itemize}
            \item 由於政策規定「利息支出不可以扣抵」所得稅，這意味著借款者向銀行借錢時，所面對的實質借款利率仍然是原來的市場利率 \(r\)。
            \item 因為借款的代價（機會成本）完全沒有改變，借款者所面對的預算限制線在借款區段（\(C_1 > Y_1\)）完全沒有移動。
            \item 結論： 既然面對的條件沒變，借款者的最適跨期消費選擇也不會改變，因此借款者的情況不變。
        \end{itemize}
        \item 對「儲蓄者」的影響：
        \begin{itemize}
            \item 由於政策規定「利息所得要課稅」，假設稅率為 \(t\)，儲蓄者面對的稅後實質存款利率下降為 \(r(1-t)\)。這會改變第一期消費與第二期消費的相對價格。
            \item 替代效果：稅後利率下降，代表「現在消費」的機會成本變低（把錢存起來賺到的利息變少了）。因此，儲蓄者會想要增加現在的消費，減少儲蓄。
            \item 所得效果：稅後利率下降，代表儲蓄者未來的利息收入變少，整體實質購買力下降。假設正常財的前提下，變窮會使人減少現在的消費，進而增加儲蓄。
            \item 結論：由於替代效果（減少儲蓄）與所得效果（增加儲蓄）的方向相反，因此儲蓄者最終儲蓄數量的增減，必須視這兩種效果的相對大小而定。
        \end{itemize}
    \end{enumerate}

    \item 考慮一簡單之線型所得稅如下：\(T = c + tI\)；\(T\) 為稅收，\(c\) 為常數，\(t\) 代表稅率、為介於零與 1 間之常數，\(I\) 為所得。以平均稅率判斷，當 \(c < 0\) 時，稅率結構為：
    \begin{enumerate}[label=\Alph*.]
        \item 累退
        \item 累進
        \item 比例
        \item 先累進後累退
    \end{enumerate}
    解答：
    \begin{itemize}
        \item Suppose $c$ is negative, and $T = c+tI$
        \[
        ATR = \frac{T}{I} = \frac{c + tI}{I} = \frac{c}{I} + t
        \]
        Since $c$ is negative, as $I$ increases $c,t$ are constant, so $\frac{c}{I}\ \text{increases} \implies \frac{c}{I} + t$ increases $\rightarrow ATR \uparrow \implies \text{Progressive}$
        \item Suppose $c$ is positive, and $T = c+tI$
        \[
        ATR = \frac{T}{I} = \frac{c + tI}{I} = \frac{c}{I} + t
        \]
        Since $c$ is positive, as $I$ increases $c,t$ are constant, so $\frac{c}{I}\ \text{decreases} \implies \frac{c}{I} + t$ decreases $\rightarrow ATR \uparrow \implies \text{Regressive}$
    \end{itemize}

    \item 若政府以均一稅制課徵所得稅，稅負為 \(T = t(Y - \bar{Y})\)，\(t\) 為稅率，\(\bar{Y}\) 為免稅額，\(Y\) 為稅前所得。若政府將 \(t\) 提高 1\%，某甲稅前所得在 \(t\) 提高前後皆固定且大於 \(\bar{Y}\)，則對某甲會造成何種影響？
    \begin{enumerate}[label=\Alph*.]
        \item 邊際稅率提高 1\%，平均稅率提高低於 1\%
        \item 邊際稅率提高 1\%，平均稅率提高超過 1\%
        \item 邊際稅率提高超過 1\%，平均稅率提高低於 1\%
        \item 邊際稅率提高低於 1\%，平均稅率保持不變
    \end{enumerate}
    解答：
    \begin{enumerate}[1.]
        \item 邊際稅率 (MTR) 的變化：
        \begin{itemize}
            \item 邊際稅率定義為增加一單位所得所增加的應納稅額，即對函數作微分：\(MTR = \frac{dT}{dY} = t\)。
            \item 當政府將稅率 \(t\) 提高 1\%（即 \(t\) 變成 \(t + 1\%\)），新的邊際稅率變為 \(t + 1\%\)。
            \item 結論：邊際稅率剛好提高了 1\%。
        \end{itemize}
        \item 平均稅率 (ATR) 的變化：
        \begin{itemize}
            \item 平均稅率定義為總應納稅額佔總所得的比例：\(ATR = \frac{T}{Y} = \frac{t(Y - \bar{Y})}{Y} = t \cdot (1 - \frac{\bar{Y}}{Y})\)。
            \item 當稅率 \(t\) 提高 1% 時，新的平均稅率為 \(ATR_{new} = (t + 1\%) \cdot (1 - \frac{\bar{Y}}{Y})\)。
            \item 平均稅率的增加幅度為：
            \[
            \Delta ATR = ATR_{new} - ATR = (t + 1\%) \cdot (1 - \frac{\bar{Y}}{Y}) - t \cdot (1 - \frac{\bar{Y}}{Y})
            \]
            \[
            \Delta ATR = 1\% \cdot (1 - \frac{\bar{Y}}{Y})
            \]
            \item 結論：甲的所得大於免稅額（\(Y > \bar{Y} > 0\)），所以分數 \(\frac{\bar{Y}}{Y}\) 會介於 0 與 1 之間，所以括號內的乘數 \((1 - \frac{\bar{Y}}{Y})\) 小於 1。
        \end{itemize}
    \end{enumerate}

    \item 下列關於拉法曲線（Laffer curve）的敘述何者正確？
    \begin{enumerate}[label=\Alph*.]
        \item 若勞動供給愈無彈性，拉法曲線中對應最高稅收的稅率愈高，故降低稅率愈有可能減少稅收
        \item 拉法曲線的最高點剛好對應後彎的勞動供給曲線折彎的那一點
        \item 因為稅收等於稅率乘以稅基，所以拉法曲線表示稅率愈高，稅收愈高
        \item 減稅必可提高勞動供給，納稅人所得因而增加，故政府稅收增加
    \end{enumerate}
    解答：
    \begin{enumerate}[label=\Alph*.]
        \item 正確：拉法曲線向下彎折，是因為高稅率會產生「替代效果（休閒變便宜，勞動減少）」，導致稅基流失。如果「勞動供給愈無彈性」，代表勞工對稅率變動的反應很小（即使稅率提高，也不會大幅減少工作時數）。此時，政府可以將稅率訂得非常高，才會開始出現稅收減少的現象。也就是說，對應最高稅收的最適稅率（$t^*$）會「愈高」。

        在這種情況下，因為最適稅率很高，目前的實際稅率很有可能落在拉法曲線的「左半邊（正斜率階段）」。在左半邊，降低稅率只會導致稅收減少，無法產生供給學派所謂「減稅帶來稅收增加」的效果。
        \item 錯誤：拉法曲線的最高點（稅收極大化）與「後彎的勞動供給曲線」的折彎點（勞動時數極大化）沒有必然的對應關係。

        勞動供給曲線後彎是因為工資上升到一定程度後，「所得效果（想多買休閒）」大於「替代效果」。而拉法曲線的最高點則是取決於整個經濟體總稅基（總勞動所得）對稅率的彈性，兩者的經濟意義與數學極值條件不同。
        \item 錯誤：拉法曲線正是用來反駁「稅率愈高，稅收愈高」的觀念。當稅率超過最適稅率（$t^*$）後，由於工作意願大幅降低導致稅基萎縮的速度大於稅率提升的速度，政府的總稅收反而會下降。
        \item 錯誤：減稅是否能提高政府稅收，取決於目前的稅率是落在拉法曲線的「左半邊」還是「右半邊」。只有當目前稅率處於極高水準（落在拉法曲線右半邊的負斜率階段，即 $t > t^*$）時，減稅才能刺激勞動供給大幅增加，進而使總稅收增加。
    \end{enumerate}

    \item 假設勞動市場為完全競爭，供給曲線為 \(W = L\)，其中，\(W\) 為工資率，\(L\) 為勞動量，需求曲線為 \(W = 600 - L\)。若政府對勞動所得課徵 50\% 之稅率，則勞動市場之稅後均衡勞動量為：
    \begin{enumerate}[label=\Alph*.]
        \item 100
        \item 200
        \item 300
        \item 400
    \end{enumerate}
    解答：略。

    \item 1.新開徵能源稅收來融通新增長期照護的支出\ 2.提高營業稅稅收來彌補所得稅的減徵。以上兩種政府租稅措施的歸宿分析方法分別為：
    \begin{enumerate}[label=\Alph*.]
        \item 1.絕對租稅歸宿，2.平衡預算歸宿
        \item 1.平衡預算歸宿，2.差異租稅歸宿
        \item 1.差異租稅歸宿，2.絕對租稅歸宿
        \item 1.差異租稅歸宿，2.平衡預算歸宿
    \end{enumerate}
    解答：
    \begin{enumerate}[1.]
        \item 情境 1：「新開徵能源稅收來融通新增長期照護的支出」對應於「平衡預算歸宿（Balanced-budget Incidence）」
        \begin{itemize}
            \item 定義：探討政府在「增加某項租稅」的同時，將這筆新增的稅收「全數用於某項特定的政府支出」，從而分析這兩項政策結合後對社會所得分配的淨影響。
            \item 分析：題目中增收「能源稅」並將其用於「長期照護支出」，這種將稅收與特定支出掛鉤、維持預算平衡的分析，是平衡預算歸宿。
        \end{itemize}
        \item 情境 2：「提高營業稅稅收來彌補所得稅的減徵」對應於「差異租稅歸宿（Differential Tax Incidence）」
        \begin{itemize}
            \item 定義： 在政府「總支出不變」且「總稅收不變」的前提下，以「增加一種租稅」來替代「減少另一種租稅」，藉此比較兩種不同稅制對所得分配的相對差異。
            \item 題目中「提高營業稅」是為了彌補「所得稅減徵」，代表政府整體的稅收規模沒有改變，為稅制結構的替換（以營業稅替代所得稅），是差異租稅歸宿。
        \end{itemize}
    \end{enumerate}
    絕對租稅歸宿 (Absolute Tax Incidence) 則是單純只探討某項租稅的增減，完全不考慮政府支出變化或其他稅制替代的問題，通常會導致總體經濟的緊縮或擴張，因此在個體經濟的租稅分析中較少被使用。

    \item 考慮比例薪資所得稅對勞動供給的影響，在什麼條件下，此租稅才不會減少勞動供給？
    \begin{enumerate}[label=\Alph*.]
        \item 休閒為正常財，且所得效果大於替代效果
        \item 休閒為正常財，且所得效果小於替代效果
        \item 休閒為劣等財，且所得效果等於替代效果
        \item 休閒為劣等財，且所得效果小於替代效果
    \end{enumerate}
    解答：
    課徵比例薪資所得稅時，勞動者的稅後工資，也就是淨工資會下降：
    \[
    w_n = (1-t)w
    \]
    當有課稅時：
    \[
    w_n \downarrow
    \]
    若休閒為正常財，淨工資下降會產生兩種效果：
    \begin{itemize}
        \item 替代效果：淨工資下降，使休閒的機會成本下降，因此勞動者會增加休閒、減少工作時數。
        \item 所得效果：淨工資下降，使勞動者實質所得減少。因休閒是正常財，所以勞動者會減少休閒、增加工作時數。
    \end{itemize}
    若要課徵比例薪資所得稅後，不會減少勞動供給。
    因此方向是：
    \[
    \text{所得效果} > \text{替代效果}
    \]
    使得：
    \[
    w_n \downarrow \quad \Rightarrow \quad L \uparrow
    \]
    也就是淨工資下降時，工作時數增加。

    \item 假設甲的效用函數為 \(X^{1/2} Y^{1/2}\)，其中 \(X\)、\(Y\) 財貨均由邊際成本固定的完全競爭市場所提供，兩財貨的市場價格分別為 \(P_X = 10\)、\(P_Y = 20\)，稟賦所得為 1,000 元。下列何種稅制對甲所產生的超額負擔最小？
    \begin{enumerate}[label=\Alph*.]
        \item 對 \(X\)、\(Y\) 財貨分別課徵稅率為 10\% 與 50\% 的從價稅
        \item 對 \(X\)、\(Y\) 財貨分別課徵每單位 1 元與 5 元的從量稅
        \item 對 \(X\)、\(Y\) 財貨均課徵每單位 10 元的從量稅
        \item 對 \(X\)、\(Y\) 財貨均課徵稅率為 50\% 的從價稅
    \end{enumerate}
    解答：要找不會使相對價格改變的課稅方式，才會讓甲的超額負擔最小。

    \item 兩期消費模型中，甲、乙兩人於第一期皆有相同的稟賦所得，第二期則皆無稟賦所得。假設甲第一期最適消費水準大於乙，分別考慮稟賦所得及利息收入皆課稅的比例所得稅，或對兩期消費皆課徵相同稅率的比例消費稅，下列關於兩人終生稅負（兩期租稅負擔折現值合計）之敘述，何者正確？
    \begin{enumerate}[label=\Alph*.]
        \item 比例所得稅下，兩人終生稅負相同
        \item 比例所得稅下，甲之終生稅負較大
        \item 比例消費稅下，兩人終生稅負相同
        \item 比例消費稅下，甲終生稅負較大
    \end{enumerate}
    解答：

    \item 在勞動市場為完全競爭下，當勞工的受補償勞動供給曲線為正斜率，而廠商是市場工資的接受者時，對廠商發放的工資課徵從價稅，會產生何種結果？
    \begin{enumerate}[label=\Alph*.]
        \item 當工資稅率愈高，則課稅所產生的超額負擔愈低
        \item 當受補償的勞動供給彈性愈高，則課稅所產生的超額負擔愈高
        \item 政府收到的稅收會超過福利損失
        \item 政府收到的稅收會等於福利損失
    \end{enumerate}
    解答：
    針對勞動市場課徵薪資稅所產生的超額負擔（\(EB\)）近似值為：
    \[
    EB = \frac{1}{2} \cdot \frac{\varepsilon_D \varepsilon_S}{\varepsilon_S + |\varepsilon_D|} \cdot wL \cdot t^2
    \]
    其中 \(\varepsilon_S\) 為受補償勞動供給彈性，\(\varepsilon_D\) 為勞動需求彈性，\(wL\) 為初始勞動所得（稅基），\(t\) 為稅率。

\end{enumerate}





\end{document}
